\documentclass[10pt,a4paper]{amsart}
\usepackage[margin=19mm]{geometry}
\usepackage{amsmath,amssymb,mathtools}
\usepackage[scr=rsfs,cal=euler]{mathalfa}
\usepackage{xfrac}
\usepackage[unicode,hidelinks,bookmarks=false]{hyperref}
\newcommand{\ie}{i.e.\ }
\newcommand{\cf}{cf.\ }
\newcommand{\resp}{respectively\ }
\newcommand{\Subsection}{\subsection}
\providecommand{\nobreakdash}{\nobreak\hbox{-}}
\setlength{\emergencystretch}{2em}
\multlinegap0pt
\usepackage{mathtools}
\usepackage{mathabx}
\usepackage{mathalfa}
\usepackage{enumitem}
\setcounter{tocdepth}{2}    % table of contents
\setcounter{secnumdepth}{3}
\numberwithin{equation}{section}     % Makes labeled equations easier to find.
\newcommand{\N}{\mbox{$\mathbb{N}$}}
\newcommand{\Z}{\mbox{$\mathbb{Z}$}}
\newcommand{\C}{\mbox{$\mathcal{C}$}}
\newcommand{\R}{\mbox{$\mathbb{R}$}}
\newcommand{\Res}{\mbox{$\mathcal{R}$}}
\newcommand{\T}{\mbox{$\mathbb{T}$}}
\newcommand{\U}{\mathcal{U}}
\def\AA{{\mathbb A}} \def\BB{{\mathbb B}} \def\CC{{\mathbb C}} \def\DD{{\mathbb D}}
\def\EE{{\mathbb E}} \def\FF{{\mathbb F}} \def\GG{{\mathbb G}} \def\HH{{\mathbb H}}
\def\II{{\mathbb I}} \def\JJ{{\mathbb J}} \def\KK{{\mathbb K}} \def\LL{{\mathbb L}}
\def\MM{{\mathbb M}} \def\NN{{\mathbb N}} \def\OO{{\mathbb O}} \def\PP{{\mathbb P}}
\def\QQ{{\mathbb Q}} \def\RR{{\mathbb R}} \def\SS{{\mathbb S}} \def\TT{{\mathbb T}}
\def\UU{{\mathbb U}} \def\VV{{\mathbb V}} \def\WW{{\mathbb W}} \def\XX{{\mathbb X}}
\def\YY{{\mathbb Y}} \def\ZZ{{\mathbb Z}}
\def\cA{\mathcal{A}}  \def\cG{\mathcal{G}} \def\cM{\mathcal{M}} \def\cS{\mathcal{S}}
\def\cB{\mathcal{B}}  \def\cH{\mathcal{H}} \def\cN{\mathcal{N}} \def\cT{\mathcal{T}}
\def\cC{\mathcal{C}}  \def\cI{\mathcal{I}} \def\cO{\mathcal{O}} \def\cU{\mathcal{U}}
\def\cD{\mathcal{D}}  \def\cJ{\mathcal{J}} \def\cP{\mathcal{P}} \def\cV{\mathcal{V}}
\def\cE{\mathcal{E}}  \def\cK{\mathcal{K}} \def\cQ{\mathcal{Q}} \def\cW{\mathcal{W}}
\def\cF{\mathcal{F}}  \def\cL{\mathcal{L}} \def\cR{\mathcal{R}} \def\cX{\mathcal{X}}
\def\cY{\mathcal{Y}}  \def\cZ{\mathcal{Z}}
\newtheorem*{teo*}{Theorem}
\newtheorem*{mainteo}{Main Theorem}
\newtheorem*{teoA}{Theorem A}
\newtheorem*{teoB}{Theorem B}
\newtheorem*{teoBp}{Theorem B'}
\newtheorem*{teoC}{Theorem C}
\newtheorem{teo}{Theorem}[section]
\newtheorem{conj}{Conjecture}
\newtheorem*{conj*}{Conjecture}
\newtheorem{quest}{Question}
\newtheorem{cor}[teo]{Corollary}
\newtheorem*{corD}{Corollary D}
\newtheorem*{af}{Claim}
\newtheorem{lema}[teo]{Lemma}
\newtheorem{prop}[teo]{Proposition}
\newtheorem{propied}{Property}
\newcommand{\bi}{\begin{itemize}}
\newcommand{\ei}{\end{itemize}}
\newcommand{\itm}[1]{\item [(#1)]}
\newtheorem{defi}[teo]{Definition}
\newtheorem{obs}[teo]{Remark}
\newtheorem{ej}[teo]{Example}
\newtheorem{conv}[teo]{Convention} 
\newcommand{\esbozo}{\vspace{.05in}{\sc\noindent Sketch }}
\newcommand{\demo}[1]{\vspace{.05in}{\sc\noindent Proof #1.}}
\newcommand{\dem}{\vspace{.05in}{\sc\noindent Proof. }}
\newcommand{\lqql}{\par\hfill {$\Box$} \vspace*{.05in}\mbox{}}
\newcommand{\finobs}{\par\hfill{$\diamondsuit$} \vspace*{.05in}}
\newcommand{\cla}[1]{\overline{#1}}
\newcommand{\eps}{\varepsilon}
\DeclareMathOperator{\interior}{int}
\newcommand{\f}{\varphi}
\newcommand{\sig}{\Sigma}
\newcommand{\en}{\subset}
\newcommand{\trans}{\tpitchfork}
\newcommand{\transv}{\pitchfork}
\DeclareMathOperator{\Diff}{Diff} 
\DeclareMathOperator{\PH}{PH} 
\DeclareMathOperator{\Emb}{Emb}
\DeclareMathOperator{\Tang}{Tang}
\DeclareMathOperator{\diametro}{diam}
\DeclareMathOperator{\Image}{Image}
\newcommand{\diff}[1]{\Diff^{#1}(M)}
\newcommand{\marg}[1]{\marginpar{\tiny{#1}}} %Este hace un margen
\newcommand{\margR}[1]{\marginpar{\flushright{\tiny{#1}}}} %Este es por si sale del lado izquierdo
\newcommand{\margem}[1]{\marginpar{{\scriptsize {#1}}}}
\begin{document}
\title[Robust transitivity versus trapping regions for partially hy]{Robust transitivity versus trapping regions for partially hyperbolic diffeomorphisms}
\author{Sylvain Crovisier; Rafael Potrie}
\date{}
\maketitle
\noindent\emph{Source and licence.} Robust transitivity versus trapping regions for partially hyperbolic diffeomorphisms. Version arXiv:2607.04496v1 (CC BY 4.0). This material is distributed under the Creative Commons Attribution 4.0 International licence, http://creativecommons.org/licenses/by/4.0/, as recorded in the arXiv OAI licence metadata for this exact version and shown on the abstract page https://arxiv.org/abs/2607.04496v1. Full rights notice: licenses/arxiv-2607.04496-CC-BY-4.0.txt.\par\smallskip
\noindent\textit{Editorial scope.} Counted unit: the original \texttt{teo} environment \texttt{teo-hyp} at source lines 382--388 together with its complete proof at lines 390--423; the delivered document keeps source lines 367--424 so that every referenced label and every citation resolves inside it. Native editable source is the paper's own concatenated TeX; the AMS wrapper is editorial formatting only. No publisher class, upstream script or unrelated artwork is redistributed.
\section{Subsets without stable connections}\label{s.surface} 

In this section we deal with invariant sets which do not have \emph{stable connections}, i.e. sets which intersect each strong stable leaves in at most one point. The following is proved in~\cite{BC-Center}.
\begin{teo}\label{teo-surface}
Let $f$ be a partially hyperbolic diffeomorphism and $K$ be a compact $f$-invariant set satisfying $\cW^s(x) \cap K = \{x\}$ for every $x \in K$. Then, there exists a submanifold $S$ such that:
\begin{itemize}
\item[--] The manifold $S$ contains $K$, is tangent to $E^{c}\oplus E^u$
and is locally invariant: $f(S)\cap S$ contains a neighborhood of $K$ in $S$.
\item[--] For any diffeomorphism $g$ $C^1$-close to $f$, any $g$-invariant compact set $K_g$ in a small neighborhood of $K$ satisfies $\cW^s_g(x) \cap K_g = \{x\}$ for every $x \in K_g$. Moreover $K_g$ is contained in
a submanifold $S_g$ which is arbitrarily $C^1$-close to $S$ as $g$ gets closer to $f$.
\end{itemize}
\end{teo}

When $f$ is $C^1$-generic and $\dim E^c=1$, this can be combined with the results of \cite{PS-IHP} (see also \cite[Corollary 2.31]{CP} or \cite{CPS}).


\begin{teo}\label{teo-hyp}
There is a dense $G_\delta$ subset $\cG \subset \mathrm{PH}^1_{c=1}(M)$ (resp. $\cG \subset \mathrm{PH}^1_{c=1,\mathrm{vol}}(M)$) such that if $f \in \cG$, then any compact $f$-invariant set $K \subset M$ satisfying $\cW^s(x) \cap K = \{x\}$ for every $x \in K$ is hyperbolic.

If $E^c|_{K}$ is uniformly expanding, then $K$ is a finite set of hyperbolic saddles.
If $E^c|_{K}$ is uniformly contracting and $K$ is an unstable lamination, then $K$ is
an attractor: there exists a trapping region $U$ such that $\bigcap_n f^n(U)=K$.
\end{teo}

\begin{proof}
Let us consider a countable basis of the topology of $M$, which consists in open sets
whose diameter is smaller than the scale  $\delta_0$ of the local product structure.
Let $\cB= \{U_n\}$ be the set of finite unions of elements of the base. Given $U_n \in \cB$
and a diffeomorphism $g$, we consider $K_n(g)$ the maximal invariant set of the closure of $U_n$. Let also $\cA_n \subset \mathrm{PH}^1_{c=1}(M)$ (resp. $\cA_n \subset \mathrm{PH}^1_{c=1,\mathrm{vol}}(M)$) be the open set of partially hyperbolic diffeomorphisms for which $K_n(g)$ is hyperbolic: $K_n(g)$ decomposes as a disjoint union of (possibly empty) invariant compact sets $K_n(g)=K_n^1(g) \cup K_n^2(g)$ such that $E^c$ is contracting on $K_n^1(g)$ and expanding in $K_n^2(g)$. Then $\cO_n := \cA_n \cup (\overline{\cA_n})^c$ is open and dense. We define $\cG = \bigcap_n \cO_n$ which is $G_\delta$ and dense.

\begin{af}
For any $f \in \cG$ and $U_n \in \cB$, if the set $K_n(f)$ satisfies $\cW^u_f(x) \cap K_n(f) = \{x\}$ for every $x \in K_n(f)$ then it is hyperbolic.
\end{af}

This concludes. Indeed if $f \in \cG$ and if $K$ is an $f$-invariant compact set which has no stable connection, we consider some $U_n\in \cB$ which is a small neighborhood of $K$.
By Theorem~\ref{teo-surface}, the set $K_n(f)$ has no stable connection, hence, by the claim, is hyperbolic. One deduces that $K\subset K_n(f)$ is hyperbolic.
When $E^c|_K$ is uniformly expanded,
the set $K$ is a uniformly expanded invariant set for the induced system $(S,f)$,
hence it is a finite union of sources for $f|_S$.
When $E^c|_K$ is uniformly contracted and $K$ is an unstable lamination,
it is a hyperbolic set saturated by its unstable manifolds for the induced system $(S,f)$,
hence is an attractor; since $f$ uniformly contracts transversally to $S$, the set $K$
is also an attractor for $(M,f)$.

\smallskip

It remains to prove the claim: we will show that there are arbitrarily $C^1$-small perturbations $g$ of $f$ such that $K_n(g)$ is hyperbolic; since $f$ belongs to the open set  $\cA_n \cup (\overline{\cA_n})^c$ this implies $f \in \cA_n$ and thus $K_n(f)$ is hyperbolic.

We first explain how to perturb in the non-conservative case and then how to adapt it to the conservative setting. Using Theorem \ref{teo-surface} we know that $K_n(f)$ is contained in a locally invariant submanifold $S$ for $f$ tangent to $E^c \oplus E^u$ at $K_n(f)$. By perturbation we can assume without loss of generality that $f\in \cA_n \cup (\overline{\cA_n})^c$ is smooth. By local invariance, there exists a neighborhood $V$ of $K_n(f)$ such that $f(S)\cap V\subset S$.
We can consider a sequence of smooth submanifold $S_k$ which converges to $S$ in the $C^1$-topology (and can be seen as a graphs over $S$ with respect to the projection along $\cW^s$). For each $k$ there exists a smooth diffeomorphism $h_k$ such that $h_k\circ f(S_k)\cap V\subset S_k$. One can furthermore require that $(h_k)$ converges to the identity in the $C^1$-topology.

We then set $g_k=h_k\circ f$. The maximal invariant set $K_{n}(g_k)$ is close to $K_n(f)$ and has to be contained in $S_k$ (because points outside $S_k$ separate from $K_n(g_k)$ by forward iterations up to a
uniform distance). The diffeomorphism $g_k$ can be chosen so that inside $S_k\cap V$ all periodic orbits are hyperbolic and each invariant one-dimensional compact submanifold contains a periodic orbit. It follows that one can apply the results of \cite{PS-IHP} to deduce that the set $K_n(g_k)$ is hyperbolic, concluding the proof. 

In the conservative case, the difficulty is to show that one can choose the perturbations $g_k$ as above to be still conservative. First, we can apply~\cite{Avila} and reduce to the case where the initial diffeomorphism $f$ is conservative and smooth. Note that there is no connected component $S'$ of $S$ which is a compact submanifold without boundary such that $K_n(f)\cap S'=S'$: this would imply that $S'$ is fixed by an iterate of $f$ and normally contracted, contradicting the volume-preservation. In particular the one can choose the neighborhood $V$ to be a submanifold with boundary such that $M\setminus V$ is connected. We build as before diffeomorphisms $h_k$ but require in addition that $h_k\circ f$ preserves the volume form on $V$. By applying Moser's trick (see~\cite[Theorem 5.1.27]{KH})
we can modify $h_k$ outside $V$ so that it is volume-preserving everywhere and still satisfies that $h_k\to \text{Id}$
for the $C^1$-topology. In particular $g_k$ is conservative too.
\end{proof}


\begin{thebibliography}{99}
\bibitem[A]{Avila} A. Avila, On the regularization of conservative maps. \emph{Acta Math.} \textbf{205} (2010), 5--18.
\bibitem[BC$_2$]{BC-Center} C. Bonatti, S. Crovisier, Center manifolds for partially hyperbolic set without strong unstable connections. \emph{Journal of the IMJ} {\bf15} (2016), 785--828.
\bibitem[CPu]{CP} S. Crovisier, E. Pujals, Essential hyperbolicity and homoclinic bifurcations: a dichotomy phenomenon/mechanism for diffeomorphisms. \emph{Inventiones Mathematicae} {\bf 201} (2015), 385--517.
\bibitem[CPuS]{CPS} S. Crovisier, E. Pujals, M. Sambarino, Hyperbolicity of the extreme bundles. \emph{In preparation}.
\bibitem[KH]{KH} A. Katok, B. Hasselblatt, Introduction to the modern theory of dynamical systems. \emph{Encyclopedia Math. Appl.} \textbf{54} Cambridge University Press, Cambridge, (1995).
\bibitem[PuSa$_2$]{PS-IHP} E. Pujals, M. Sambarino, Density of hyperbolicity and tangencies in sectional dissipative regions. \emph{Ann. Inst. H. Poincar\'e C Anal. Non Lin\'eaire} {\bf 26} (2009), 1971--2000.
\end{thebibliography}
\end{document}
